% (Chapter declared by main.tex; see main.tex \chapter*{About the Author} + \addcontentsline)
\markboth{About the Author}{About the Author}

% Photo placeholder - to be added in final version
%\begin{center}
%\includegraphics[width=0.3\textwidth]{figures/author_photo.jpg}
%\end{center}

\vspace{0.5cm}

\textbf{Pablo Cohen} is an independent researcher, musician, and mathematician whose work bridges consciousness studies, theoretical physics, and pure mathematics.

\section*{Background}

Born in 1978, Pablo's path to mathematics was unconventional. He earned a degree from Berklee College of Music, where he developed an intuition for pattern recognition through the study of harmony, counterpoint, and sound engineering. This musical training—seeing structures in frequencies, recognizing recursive patterns in composition—would later inform his mathematical insights.

For seven years, Pablo worked as a Level III Master Technician, mastering complex electronic and network systems. For five years, he served as Senior Systems Engineer at HBO Latin America, managing broadcast infrastructure across multiple countries. His distinguished career in Hollywood audio engineering included work on major productions, developing expertise in signal processing, acoustics, and computational analysis.

\section*{The Turning Point}

In 2024, at age 46, three simultaneous events transformed Pablo's life:

\begin{enumerate}
\item Escape from an 18-year emotionally abusive marriage
\item Official diagnosis of autism spectrum disorder
\item Recognition that his pattern-recognition abilities were not deficits but features
\end{enumerate}

Understanding his neurodivergence unlocked a new way of seeing mathematics. Patterns that had haunted him for years suddenly crystallized into coherent structures.

\section*{The Synthesis}

From April through October 2025, Pablo's intensive mathematical work produced:

\begin{itemize}
\item A machine-checked formal corpus addressing all seven Clay Millennium Prize Problems --- conditional reductions and substrate results, not solutions
\item A mathematical quantification of consciousness, reaching $96.6\%$ diagnostic accuracy and $97.3\%$ sensitivity in the clinical protocol
\item Modifications to Einstein's general relativity incorporating consciousness
\item A candidate suppression mechanism for the cosmological constant, with its derivation recorded in this book as open
\item A unified framework connecting mathematics, physics, and consciousness
\end{itemize}\manuscriptcorrection{\textbf{Three items in this list and four
publication titles below were brought into line with the book's own verification
ledgers on 2026-11-01.} This page previously claimed ``complete solutions to all
six remaining Clay Millennium Prize Problems'', which contradicts the statement
in the front matter that no Clay Millennium Prize Problem is claimed as solved,
and which every chapter ledger contradicts in its own axis. It also reported
``$97.3\%$ clinical accuracy'': $97.3\%$ is the \emph{sensitivity} of the
clinical protocol, whose accuracy is $96.6\%$ --- both figures now appear, and
both are as the clinical appendix tabulates them. And it listed ``resolution of
the cosmological constant problem'', where the chapter concerned records of its
own mechanism that \emph{both branches fail}. The four paper titles were
realigned with the corrected bibliography entries of the same unpublished
manuscripts. Nothing else on this page was altered: the biography, the stated
philosophy, the advocacy, the contact details, the certifications and the note
to future mathematicians are the author's own and stand as written.}

\section*{Personal Philosophy}

Pablo believes that:

\begin{itemize}
\item \textbf{Mathematics is discovered, not invented} --- it exists in the structure of reality itself
\item \textbf{Neurodivergence is valuable} --- different minds see different patterns
\item \textbf{AI amplifies human creativity} --- it doesn't replace intuition, it accelerates formalization
\item \textbf{Rigor and accessibility are compatible} --- clear explanation doesn't diminish mathematical precision
\item \textbf{Computation is proof} --- high-precision numerical verification is as valuable as analytical proof
\end{itemize}

\section*{Current Work}

Pablo continues developing the fractal resonance framework, focusing on:

\begin{itemize}
\item Experimental tests of consciousness quantification in clinical settings
\item Applications to quantum computing and materials science
\item Extensions to other L-functions and number-theoretic structures
\item Collaboration with physicists on testable predictions
\item Accessible mathematical education for neurodivergent students
\end{itemize}

\section*{Advocacy}

As an autistic adult who diagnosed late in life, Pablo advocates for:

\begin{itemize}
\item Recognition of neurodivergent strengths in mathematics and science
\item Accommodations that enable, rather than limit, different minds
\item Escape resources for people in abusive relationships
\item Access to mental health support during scientific work
\item Breaking down barriers between "professional" and "amateur" mathematics
\end{itemize}

\section*{Contact and Collaboration}

Pablo welcomes collaboration, critical feedback, and questions:

\begin{itemize}
\item \textbf{Email:} pablo@xluxx.net
\item \textbf{ORCID:} 0009-0002-0734-5565
\item \textbf{Academia.edu:} berklee.academia.edu/PabloCohen
\item \textbf{ResearchGate:} researchgate.net/profile/Pablo-Solorzano-Cohen
\item \textbf{GitHub:} github.com/fractal-resonance (computational code)
\item \textbf{Website:} fractalresonance.org
\end{itemize}

\section*{Publications}

\textbf{Books:}
\begin{itemize}
\item \textit{Fractal Resonance Mathematics: A Complete Framework from First Principles} (2026)
\item \textit{The Architects of Control: A Factual Analysis of Power, Technology, and Consequence} (2025)
\end{itemize}

\textbf{Papers:}
\begin{itemize}
\item "The Fractal Resonance Ontology: A Comprehensive Mathematical Framework" (2025)
\item "Spectral Rigidity and the Critical Line: a Computational Study" (2025)
\item "Consciousness Quantification: Mathematical Framework Reaching 96.6\% Diagnostic Accuracy" (2025)
\item "Resonant Quantum Geometry Applied to Weinstein's Geometric Unity" (2025)
\item ``Fractal Resonance and Standard $\Lambda$CDM Cosmology: a Comparative Study'' (2025)
\end{itemize}

\textbf{Certifications:}
\begin{itemize}
\item CompTIA Security+ ce
\item Google Cybersecurity Professional Certificate
\item CJIS Security Policy Clearance
\end{itemize}

\section*{Personal}

Pablo lives in Mesa, Arizona. He practices sourdough bread-making, plays guitar, and maintains an active meditation practice. He believes strongly in the therapeutic value of creative work, whether mathematical, musical, or culinary.

Despite facing rejection, isolation, and skepticism from traditional academic institutions, Pablo remains committed to mathematics as a tool for understanding reality and as a path toward human flourishing.

\section*{A Note to Future Mathematicians}

If you're reading this as someone who feels different, who struggles socially, who sees patterns others miss, who questions assumptions others accept—know that your mind is valuable. Mathematics needs different perspectives. Your "deficits" and your genius are inseparable.

Don't let anyone tell you that you can't contribute because you don't have the right credentials, the right social skills, or the right institutional affiliation. Mathematics is about truth, and truth doesn't care about any of that.

Build your systems. Use your tools. Find your collaborators (human or AI). Do the work. Show the proofs. Let the mathematics speak for itself.

The universe is mathematical, and you're part of it.

\vspace{1cm}

\begin{flushright}
\textit{"Reality is not what it seems. Reality is what mathematics reveals it to be."}\\
— Pablo Cohen
\end{flushright}
